% Short running form: the full title is wider than the box the running head
% leaves beside the Sunday's name.
\runninghead{The Mind of Christ Jesus}
\properlane{obedient-son}{The Mind That Was in Christ Jesus: The Son Who Obeyed unto Death}
\label{sec:the-mind-of-christ-jesus}

Read from the second reading, the formulary turns on Christ's humble
obedience, which is both the pattern and the power of Christian humility and
charity.
He who was in the form of God emptied himself, obeyed the Father as a Son unto
death on a cross, and was exalted; the disciple is to have the same mind,
counting others better than himself, obeying against his own will, and laying
down his life for the brethren. The Collect's claim that God shows his power
most in sparing and in mercy then names the same movement from God's side:
Aquinas says that by sparing and mercy God leads men to share the infinite
good, and Schuster calls the sinner's restoration a condescension greater than
creation. That the power of the Collect is a power that stoops is the
editor's joining of the two; the stooping is Schuster's word, not Aquinas's.

\subsection{Paul's appeal: one mind and lowliness}

Paul's appeal begins with a ``therefore'' that ties it to the conflict the
Philippians share with him (1:27--30), and it asks for one mind and one charity
and for nothing done ``through contention'' or ``vain glory''. Chrysostom, on
these verses, calls vainglory ``the cause of all evil'', and he will not let
the command to count others better remain a phrase: it is fulfilled ``if thou
not only sayest it, but art fully assured of it''. ``Haughtiness is the first
act of ingratitude, for it denies the gift of grace.'' He closes the homily
with Paul's other sentence, ``Not the hearers of a law, but the doers of a law
shall be justified'', and adds, ``Yea, knowledge itself condemneth, when it is
without action and deeds of virtue'' (\work{Homilies on Philippians} 5).

Aquinas explains how the command can be kept without pretence. The unity Paul
asks for is \latin{interius in affectu, et exterius in effectu}, inward in
affection and outward in effect, and he cites John's letter for it:
\latin{Non diligamus verbo, neque lingua, sed opere, et veritate} (1~Jn 3:18).
Humility answers pride: \latin{Sicut enim pertinet ad superbiam quod homo se
extollat supra se, ita ad humilitatem quod homo se subiiciat secundum suam
mensuram}, as it belongs to pride that a man raise himself above himself, so it
belongs to humility that he subject himself according to his measure. How can
the better man honestly count the worse his superior? \latin{nullus est sic
bonus, quin in eo sit aliquis defectus, et nullus est sic malus, quin habeat
aliquid boni}; and if another is wholly bad, \latin{Si ergo illum non praeponas
ratione suae personae, praeponas ratione imaginis divinae}, if you do not
prefer him by reason of his person, prefer him by reason of the divine image
(\work{Super Philippenses} c.~2, lect.~1, pp.~375--376). The shorter form of
the reading ends here, with the call \latin{Hoc enim sentite in vobis, quod et
in Christo Iesu}, ``let this mind be in you, which was also in Christ Jesus''.
Where that form is read, only the exhortation and the name of Christ's mind
remain, which is what the \work{Ordo} sets over the reading. Christ's
self-emptying, his obedience unto death and his exaltation, of which the
witnesses below speak, belong to the longer form alone.

\subsection{He emptied himself}

Chrysostom begins his next homily by noticing Paul's method: ``in exhorting
them to humility, he brought forward Christ''. The reason is that ``Nothing
rouses a great and philosophic soul to the performance of good works, so much
as learning that in this it is likened to God.'' He reads the hymn as an answer
to the appeal before it: ``It is written, `He emptied Himself' in reference to
the text, `each counting other better than himself.' '' Christ's self-emptying
is the model of that counting, and it was free, ``had He not chosen it of His
own accord, and of His own free will, it would not have been an act of
humility'' (\work{Homilies on Philippians} 6--7). Aquinas says the same in a
sentence: \latin{hic hortatur ad virtutem humilitatis, exemplo Christi}, here
Paul exhorts to the virtue of humility by the example of Christ.

Both insist that the one who emptied himself did not cease to be what he was.
The equality with God, Chrysostom says, Christ had ``not by seizure, but as his
own by nature'' (\work{Homilies on Philippians} 7); and Aquinas, on \latin{semetipsum exinanivit}: he did not empty
himself of the divinity, \latin{quia quod erat, permansit: et quod non erat,
assumpsit}, for what he was, he remained, and what he was not, he took
(\work{Super Philippenses} c.~2, lect.~2, p.~377). The example has force because
the one who gives it is God. Augustine puts that in a single sentence while
expounding the washing of the feet: ``so great is the beneficence of human
humility, that even the Divine Majesty was pleased to commend it by His own
example; for proud man would have perished eternally, had he not been found by
the lowly God.'' And he draws the consequence: ``as he was lost by imitating
the pride of the deceiver, let him now, when found, imitate the Redeemer's
humility'' (\work{Tractates on John} 55.7).

\subsection{Obedient unto death}

``He humbled himself, becoming obedient unto death.'' Chrysostom answers those
who took the obedience for a sign that the Son was less than the Father: ``He
became obedient as a Son to His Father; He fell not thus into a servile state,
but by this very act above all others guarded his wondrous Sonship, by thus
greatly honoring the Father'' (\work{Homilies on Philippians} 7). Aquinas
defines the virtue: \latin{Modus humiliationis, et signum humilitatis est
obedientia: quia proprium superborum est sequi propriam voluntatem}, obedience
is the manner of humiliation and the sign of humility, because it is proper to
the proud to follow their own will. It is set in the Passion \latin{quia prima
praevaricatio est facta per inobedientiam}, because the first transgression was
made by disobedience; \latin{obedientia dat meritum passionibus nostris},
obedience gives merit to our sufferings; and \latin{tunc est obedientia magna,
quando sequitur imperium alterius contra motum proprium}, obedience is great
when it follows another's command against one's own inclination. Christ refused
neither death nor death's most shameful kind (\work{Super Philippenses} c.~2,
lect.~2, p.~379).

Augustine dwells on the last words of the verse. Mercy comes first, then
judgment, and the first mercy is that ``The Creator of man deigned to become
man''. To this he added, step by step, that he was rejected, dishonoured and
put to death; ``but even this was not enough, it was by the death of the
cross. For when the apostle was commending to us His obedience even unto death,
it was not enough for him to say, `He became obedient unto death;' \ldots\ but
he added, `even the death of the cross' '' (\work{Tractates on John} 36.4).

The exaltation follows. ``He was obedient to the uttermost, wherefore He
received the honor which is on high'', Chrysostom says, and draws the lesson:
``Let us too not suppose then that we descend from what is our due, when we
humble ourselves.'' He also names the failure the hymn exposes: ``when we
glorify Him rightly, but live not rightly, then do we especially insult Him''
(\work{Homilies on Philippians} 7). Aquinas reads vv.~9--11 as the reward,
\latin{Hic commendat eius praemium: quod est exaltatio, et gloria}, and cites the
Lord's saying that he who humbles himself shall be exalted (Lk 14:11; 18:14;
\work{Super Philippenses} c.~2, lect.~3, p.~380).

\subsection{The son who said no and went}

No witness of this interpretation joins Philippians 2 to the parable of the two
sons or to Ezekiel; the join that follows is the editor's. Set beside the
hymn, the parable shows two sons who both fall short of the Son. One said yes
and did not go. The other said no, and then obeyed against his own first word,
which is what Aquinas calls great obedience, \latin{contra motum proprium}; but
Aquinas says it of Christ in Philippians, not of the son in the parable. The
Son of the hymn is the one whose yes was done to the end, and whose obedience
the first son's belated going imitates from far off.

The first reading is heard here through Augustine's reading of the complaint
it answers. Those who say ``The way of the Lord is not right'' are the froward,
and what seems froward to them is ``that Thou wilt make them whole that confess
their sins''; God humbles those who are ``ignorant of God's righteousness, and
seek to establish their own'' (\work{Enarrationes in Psalmos} 17.28). The
Lord's way heals the humble who confess and humbles the proud who claim their
own justice. Augustine says this of the complaint, not of Paul's hymn, and the
likeness between that way and the path Christ walked in Philippians is the
editor's.

The psalm prays for the same disposition. The Lord ``will guide the mild in
judgment: he will teach the meek his ways'', and Augustine says who these are:
``those that follow His will, and do not, in withstanding It, prefer their
own''; the Lord teaches his ways ``not to those that desire to run before, as
if they were better able to rule themselves'' (\work{Enarrationes in Psalmos}
24.9). Bellarmine adds that the proud are not excluded, since grace first
humbles them and then guides them (on Ps~24:9, quoted under the first
interpretation). The acclamation names the disciple's part: ``they follow me''.
Aquinas counts following as the third of four things in the verse,
\latin{Tertium, quod est ex parte nostra, est imitatio ad Christum}, the third,
which is on our part, is the imitation of Christ, and cites Peter: ``Christ
also suffered for us, leaving you an example that you should follow his
steps'' (1~Pet 2:21; \work{Super Ioannem} c.~10, lect.~5, p.~668).

\subsection{The life laid down, at the table}

The second Communion antiphon states what this obedience means for the
disciple: ``he hath laid down his life for us: and we ought to lay down our
lives for the brethren''. Augustine calls the verse the ``perfection of
love'', and he will not let the perfection discourage the beginner: ``Do not
too soon despair of thyself. Haply, it is born and is not yet perfect; nourish
it.'' Its beginning is in the verse that follows, the brother in need: ``If thou
art not yet equal to the dying for thy brother, be thou even now equal to the
giving of thy means to thy brother'' (\work{Homilies on the First Epistle of
John} 5.11--12). In his tractates on the Gospel he brings the verse to the
altar. Reading the proverb about sitting at a ruler's table, he asks, ``For what
is the table of the ruler, but that from which we take the body and blood of
Him who laid down His life for us? And what is it to sit thereat, but to
approach in humility?'' To put forth the hand there, knowing that one must
make the like preparation, is ``that, as Christ laid down His life for us, so
we also ought to lay down our lives for the brethren''; the martyrs did so, and
are commemorated at that table so that they may pray for us (\work{Tractates on
John} 84.1). He adds at once that no martyr's blood is shed for the remission of
sins as Christ's was (84.2). Aquinas makes the same verse the pattern of the
pastor: the good shepherd's office is charity, and of the shepherd who bears the
loss of bodily life for the flock, \latin{Huius autem doctrinae Christus nobis
exemplum praebuit}, citing 1~John 3:16 (\work{Super Ioannem} c.~10, lect.~3,
p.~660). Gregory the Great, on the shepherd who lays down his life, has him say
\latin{ea charitate qua pro ovibus morior quantum Patrem diligam ostendo}: by
the charity with which I die for the sheep I show how much I love the Father
(\work{Homiliae in Evangelia} 14.4). In Gregory's sentence the Son's love for
the Father and his death for the sheep are one act; that this love is the
obedience of Philippians 2:8 is the editor's joining.

The first Communion antiphon is the prayer of the lowly. On \latin{in
humilitate mea} Augustine says that hope \latin{data est humilibus}, is given
to the humble, and cites the Lord's saying that he who humbles himself shall be
exalted. He understands the lowliness first as humiliation suffered,
\latin{qua quisque humiliatur aliqua tribulatione vel deiectione, quam meruit
eius superbia}, by which someone is humbled by some tribulation or casting
down that his pride deserved, and not only the humility of one who confesses
his sins (\work{Enarrationes in Psalmos} 118, s.~15.2, PL 37, 1541). Hilary
likewise reads it as suffering endured, \latin{cum contemnitur, cum inridetur,
cum iniuriis uexatur, cum contumeliis inhonestatur}, when one is despised,
mocked, harassed with injuries and dishonoured with insults, and says that the
servant knows \latin{humilitatem suam glorioso esse honoris praemio
munerandam}, that his humiliation is to be rewarded with a glorious prize of
honour (\work{Tractatus super Psalmos} 118, Zain~2, CSEL 22, p.~419). Where the
longer form is read and the first antiphon is chosen, the servant's
\latin{humilitate} and the hymn's
\latin{humiliavit semetipsum} stand in the same Mass, and humiliation followed
by glory is the shape of both; that is an observation about the words, and no
witness makes it.

\subsection{The Missal's prayers in this reading}

The Collect names the power at work. Aquinas's second reason why omnipotence is
shown most in sparing is that \latin{parcendo hominibus et miserando, perducit
eos ad participationem infiniti boni}, by sparing and showing mercy God leads
men to share in the infinite good (\work{Summa theologiae} I, q.~25, a.~3,
ad~3); Schuster's is that God goes down into the abyss that separates him from
evil to draw the sinner out, in words quoted under ``Each Element in Its
Setting'' (\work{The Sacramentary}, vol.~3, p.~121). The petition asks
that those who run toward God's promises share the goods of heaven, and that
is the share in the infinite good of Aquinas's reply. Joining this to the
hymn's descent and exaltation is the editor's.

The Entrance Antiphon asks God to ``give glory to thy name'', and the
Missal's \latin{da gloriam nomini tuo} is the Clementine's at Daniel 3:43. In
the longer form of the second reading God gives Jesus ``a name which is above
all names''. The shared word is in the texts; the join is the editor's. The
Prayer over the Offerings asks that the assembly's offering be accepted; in
this interpretation the offering is taken into the self-offering of the Son who
obeyed unto death. The Prayer after Communion asks that those who share the
Lord's sufferings in proclaiming his death be his fellow heirs in glory: where
the longer form is read, that is humiliation unto death and then exaltation,
now asked for the communicants. The prayer's \latin{coheredes} and
\latin{compatimur} stand together in Romans 8:17; that too is an observation
about the words.

\subsection{Where this interpretation strains}

No witness of this interpretation joins Philippians 2 to the parable or to
Ezekiel, and the second reading follows its own course through the letter. The
interpretation's hold on the whole formulary therefore rests on the Collect
and the Communion antiphons, where the witnesses do speak, and elsewhere on the
editor's joining. Chrysostom's homilies on 2:5--11 are chiefly doctrinal,
against those who made the Son less than the Father, and his moral application
is brief. Christ is not only an example: Chrysostom and Aquinas both insist
that the one who emptied himself kept the form of God, and a reading that made
the hymn a lesson in imitation alone would misstate them. Where the shorter form
is read, this interpretation keeps only Paul's appeal and the name of Christ's
mind.

\subsection{The four senses of this interpretation}

\begin{fourSenses}
\item[Literal.] Paul asks the Philippians for one mind and humility, each
counting the others better than himself, and sets before them the mind of
Christ Jesus; in the longer form, the Christ who, being in the form of God,
emptied himself, obeyed unto the death of the cross, and was given the name
above all names.

\item[Allegorical.] In the longer form, Christ is the Son who obeyed the Father
and so honoured him (Chrysostom), and the lowly God who found proud man
(Augustine). Where the second Communion antiphon is chosen, he is also the
shepherd who lays down his life, whose body and blood are taken at the ruler's
table (Augustine).

\item[Moral.] Count others better than yourself, by the image of God in them if
by nothing else, and obey against your own inclination (Aquinas); where the
second antiphon is chosen, begin to lay down your life by giving of your means
to a brother in need (Augustine); do not glorify Christ rightly and live
wrongly (Chrysostom).

\item[Anagogical.] Exaltation follows humility (Hilary, on the servant's
lowliness, where the first antiphon is chosen; Aquinas, on Philippians, where
the longer form is read); the communicants ask to be fellow heirs in glory with
him whose death they proclaim; where the longer form is read, every knee shall
bow at the name of Jesus.
\end{fourSenses}
